\chapter{From diagnosis to a research question}\label{ch:research}
\question{What would an additional physical-action account need to provide?}

Suppose the clinic's community is considering a new pump. The financial question asks whether it can be purchased and maintained within a budget. The engineering question asks whether it can deliver the required service. The access question asks who benefits and who may still be excluded. The environmental question asks what the whole represented process draws from and leaves in its surroundings.

A useful new account must improve our ability to answer some of these questions without pretending to answer all of them. It must also be worth the resources required to measure, verify and govern it. More records are not automatically better records.

\section{There is already work to build on}
Environmental accounting is not an empty field awaiting EBU. The UN's environmental-economic accounting system includes a framework for ecosystem extent, condition and services. Its physical accounts connect information about ecosystems with human activity. The framework was adopted in March 2021; that statement does not mean every possible monetary valuation has identical statistical-standard status.\cite{seea}

Environmental regulation, public provision and engineering control also already connect observations to action. Monetary policy has its own objectives and tools, as Chapter~\ref{ch:money} showed. Their purposes, spatial scales and information requirements differ. Their existence is a reason to make the proposed contribution precise.

The EBU question considered here concerns a declared action, a permitted local description, a finite before-and-after comparison and an auditable signed result. A possible contribution would lie in how those elements work together in a particular domain. This is a research proposal, not a claim of priority over every earlier accounting or allocation method.

\clearpage
\section{Seven questions for a proposed account}
Before calculating a pump-related value, we would ask:

\begin{enumerate}
\item Which physical quantities are represented, and where is the boundary?
\item Which states and actions are physically feasible?
\item What reference and scale give a change its declared meaning?
\item What information is allowed at the moment of valuation?
\item Does the calculation include the complete finite change, including shared effects?
\item Which burdens are already in the state account, and which justified residual process burdens remain?
\item What evidence connects the proposed action to what actually occurred?
\end{enumerate}

These are design requirements, not promises of success. A model can answer them clearly and still be a poor representation of the domain. Clear failure is valuable: it tells us where a claim stops instead of letting an attractive equation conceal a missing physical quantity.

For example, an account of water volume cannot by itself certify water quality. If quality is decisive, the modeller must represent it appropriately or report the account's limitation. Declaring the volume calculation exact does not repair that omission. The same reasoning applies to infrastructure condition, delayed damage and access.

\section{What would count against the proposal?}
The proposed architecture would face a mathematical failure if a supposedly exact local value omitted an affected term. It would face a measurement problem if accepted quantities could not be connected reliably to observations. It would face a practical problem if verification cost exceeded the benefit of the information. It would face an institutional problem if apparently neutral rules systematically denied essential access without a defensible justification.

These failures need different investigations. An algebraic proof can settle the first kind of claim under explicit assumptions. It cannot settle all the others. A software check may show that a formula was implemented faithfully while leaving the formula's real-world relevance unresolved.

The prospective Gaussian research programme asks a more restricted behavioural question: what happens when actors choose among feasible actions under a separately declared capacity rule? A useful study must permit adverse answers such as trapping, repeated refusal or poor recovery. The motivation for the project cannot be used to classify every outcome as success.

\section{Commitments and choices}
Transparency, contestability and the protection of essential access are commitments that guide the inquiry. They do not emerge from differentiating a potential. Who declares a reference, who can challenge a measurement and who carries the cost of a necessary action remain institutional questions.

Consider repair work that uses materials now but may improve service later. The current physical effects should be recorded under the current account. Evidence of changed capacity can justify a prospective update to the model. Responsibility for financing the repair and compensation for the worker need their own arrangements. A forecast of future usefulness should not be inserted as an immediate observed improvement.

This way of framing the project is demanding. It asks EBU to be clear enough to audit and modest enough to reject when it fails. It also leaves room for a useful result short of a complete economic system: a particular accounting method might help in one domain and fail in another.

\practice{A pump repair leaves the tank level unchanged. Has the repair achieved nothing?}{The stock measurement alone cannot answer. A repair may change equipment condition or future feasible delivery. Those require their own evidence and representation. The immediate material use can still be recorded. Missing a capacity coordinate is a model limitation, not proof that the work was worthless.}

\carry{To ask an account to preserve function, we must first explain what preserving function means. The next chapters introduce homeostasis and balance without assuming that a valuation formula already supplies a controller.}
